\subsection{环的维数与理想的高度}

定义 5. —（交换）环 A 的克鲁尔维数，或简称维数，是拓扑空间  的克鲁尔维数，记为  （II, § 4, no. 3, definition 4）。如果  是 A 的素理想，则 A 在 \_ 处的维数记为 ，即 \_。

映射 \(\mathfrak{p} \mapsto \mathrm{V}(\mathfrak{p})\) 是从 A 的素理想集合到 \(\operatorname{Spec}(A)\) 的不可约闭子集集合的一个反序双射（同上，命题 14 的推论 2）。因此，A 的维数是 A 的素理想链的长度集合的上确界；它等于 \(-\infty\)、\(+\infty\) 或一个正整数。

设 \(\mathfrak{p} \in \operatorname{Spec}(\mathbf{A})\)；集合 \(\operatorname{Spec}(\mathbf{A})_f\)（其中 \(f\) 遍历 A）构成 \(\operatorname{Spec}(\mathbf{A})\) 的拓扑基，并且 \(\mathfrak{p}\) 属于开集 \(\operatorname{Spec}(\mathbf{A})_f\) 当且仅当 \(f\) 不属于 \(\mathfrak{p}\)。因此，\(\dim_{\mathfrak{p}}(\mathbf{A})\) 是数 \(\dim(\mathbf{A}_f)\) 的下确界，其中 \(f\) 遍历 \(\mathfrak{p}\)（II, § 5, no. 1, 命题 1）。

例。--- 1) 当且仅当 A 退化为 0 时，我们有  。要使 \textless 成立，必要且充分条件是 A 的每个素理想都是极大理想。维数为 0 的整环就是域。当且仅当它是阿廷环时，诺特环的维数为 {}（IV, § 2, no. 5, prop. 9）。

\begin{enumerate}
\def\labelenumi{\arabic{enumi})}
\setcounter{enumi}{1}
\item
Dedekind 环是维数 \(\leqslant 1\) 的诺特整闭环（VII, § 2, no. 2, th. 1）。更一般地，根据 V, § 1, no. 2, 命题 9 的推论 2，一个环是 Dedekind 环的有限直积，当且仅当它是诺特的、约化的、在其全分式环中整闭的且维数 \(\leqslant 1\)。
\item
如果 A 是一个赋值环（VI, § 1, No. 2, 定义 2），则其维数等于该赋值的高度（VI, § 4, No. 4, 命题 5）。
\item
设 A 是一个环。我们有
\end{enumerate}

\[
\dim (\mathrm{A} [ \mathrm{X} ]) \geqslant \dim (\mathrm{A}) + 1.
\]

事实上，如果 \(\mathfrak{p}_0 \subset \ldots \subset \mathfrak{p}_n\) 是 A 的素理想链，长度为 \(n\)，则我们得到 A[X] 的素理想链 \(\mathfrak{p}_0' \subset \ldots \subset \mathfrak{p}_{n+1}'\){}{}，其长度为 \(n+1\)，其中规定 \(\mathfrak{p}_i' = \mathfrak{p}_iA[X]\)（\(0 \leqslant i \leqslant n\)），以及 \(\mathfrak{p}_{n+1}' = \mathfrak{p}_nA[X] + XA[X]\)。

同样的论证给出不等式 \(\dim(\mathrm{A}[[\mathrm{X}]]) \geqslant \dim(\mathrm{A}) + 1\)。由此归纳得到不等式

\[
\begin{array}{l} \dim (\mathrm{A} [ \mathrm{X} _ {1},..., \mathrm{X} _ {n} ]) \geqslant \dim (\mathrm{A}) + n, \\ \dim (\mathrm{A} [ [ \mathrm{X} _ {1},..., \mathrm{X} _ {n} ] ]) \geqslant \dim (\mathrm{A}) + n. \end{array}
\]

稍后（§ 3, no. 4, 命题 7 的推论 3 以及命题 8 的推论 3）我们将证明，当 A 为诺特环时，前面两个公式中的等号成立。

* 5) 设 X 是一个复解析流形。如果 X 在 X 的一点 x 处具有复维数 n，则 X 上 x 处解析函数芽的局部环的维数为 n。*

\begin{enumerate}
\def\labelenumi{\arabic{enumi})}
\setcounter{enumi}{5}
\item
设 \(k\) 是一个域，A 是一个非零整 \(k\)-代数。则有 \(\dim(A) = 0\)。这由 V, § 2, no. 1 的命题 1 的推论以及 \(\dim(k) = 0\) 这一事实得出。
\item
如果 n 是 A 的幂零理想，则 Spec(A) 与 Spec(A/n) 同胚（II, § 4, no. 3, 注记）。因此有 dim(A/n) = dim(A)；特别地，如果 A\(_{red}\) 是环 A 除以其幂零根得到的商环，则有 dim(A) = dim(A\(_{red}\))。
\item
存在无限维的诺特环（p. 83, 练习 13）。我们将在下文（§ 3, no. 1, 命题 2 的推论 1）看到，每个诺特局部环都有有限维数。
\end{enumerate}

命题 6. --- 设 A 是一个环。

\begin{enumerate}
\def\labelenumi{\alph{enumi})}
\item
  如果 \(\mathfrak{a}\) 是 A 的一个理想，则有 \(\dim (\mathrm{A} / \mathfrak{a})\leqslant \dim (\mathrm{A}).\)
\item
  如果 S 是 A 的一个乘法子集，则有 \(\dim (\mathbf{S}^{-1}\mathbf{A})\leqslant \dim (\mathbf{A}).\)
\item
有 \(\dim(A) = \sup_{\mathfrak{p}} \dim(A/\mathfrak{p})\)，其中 \(\mathfrak{p}\) 遍历 A 的极小素理想集合。
\item
如果 A 只有有限个极小素理想（例如 A 为诺特环时就是如此（II, § 4, no. 3, 命题 14 的推论 3）），且 \(\mathfrak{p}\) 是 A 的素理想，则有
\end{enumerate}

\[
\dim_ {\mathfrak {p}} (\mathrm{A}) = \sup _ {\mathfrak {q}} \dim_ {\mathfrak {p} / \mathfrak {q}} (\mathrm{A} / \mathfrak {q}),
\]

其中 \(\mathfrak{q}\) 遍历包含于 \(\mathfrak{p}\) 中的 A 的极小素理想集合。

\begin{enumerate}
\def\labelenumi{\alph{enumi})}
\setcounter{enumi}{4}
\tightlist
\item
设  是 A 的一个不包含于任何极小素理想的理想；则有 。特别地，如果 A 是整环，则对 A 的每个非零理想  都有 。
\end{enumerate}

根据 II, § 4, no. 3 中的注记，如果  是 A 的一个理想，则拓扑空间  同胚于  的闭子空间。断言  由此及 no. 1 的命题 1 a) 得出。断言  由同上命题 13 的推论得出。根据同上命题 14 的推论 2， 的不可约分支同胚于空间 ，其中  是 A 的一个极小素理想；断言  由 no. 1 的命题 1 的推论得出。在  的假设下，空间  只有有限个不可约分支；包含  的  的不可约分支是集合 ，其中  是包含于  中的极小素理想。断言  于是由 no. 1 的命题 1 的推论 b) 得出。

最后证明 \(e\)。我们要证明：对于 A 的任意素理想链 \(\mathfrak{p}_{0} \subset \ldots \subset \mathfrak{p}_{n}\)，只要 \(\mathfrak{a} \subset \mathfrak{p}_{0}\)，就有 \(\dim(A) \geqslant n + 1\)。根据关于 \(\mathfrak{a}\) 的假设，存在 A 的一个素理想 \(\mathfrak{p}_{-1}\)，它包含于 \(\mathfrak{p}_{0}\) 且不同于 \(\mathfrak{p}_{0}\)，从而 \(\mathfrak{p}_{-1} \subset \mathfrak{p}_{0} \subset \ldots \subset \mathfrak{p}_{n}\) 是 A 的一条长度为 \(n + 1\) 的素理想链。

注 1. --- 设 \(\rho: A \to B\) 是一个环同态。则  是数  的上确界，其中 \(\rho(p)\) 遍历 A 的极小素理想；事实上，对于 B 的每个极小素理想 \(p\)，存在 A 的一个极小素理想 \(q\)，使其包含于 \(\mathfrak{p}\)，并满足 \(\rho^{-1}(\mathfrak{q})\)（II, § 2, no. 6, 引理 2），并且有

\[
\dim (\mathbf {B} / \mathfrak {q}) \leqslant \dim (\mathbf {B} / \rho (\mathfrak {p})\mathbf {B}) \leqslant \dim (\mathbf {B})
\]

根据命题 6 a)，由此再根据命题 6 c) 得到结论。

定义 6. --- 环 A 的理想  在  中的余维数称为理想  的高度，记为 。

假设 A 是整环。那么，按照 VII, § 1, no. 6 的定义 4，A 中高度为 1 的素理想，正是按上述定义高度为 1 的素理想。

命题 7. --- a) A 的素理想  的高度，是满足 \(\mathfrak{p}_{0} \subset \ldots \subset \mathfrak{p}_{n}\) 且 \(\mathfrak{p}_{n} = \mathfrak{p}\) 的 A 的素理想链长度的上确界。

\begin{enumerate}
\def\labelenumi{\alph{enumi})}
\setcounter{enumi}{1}
\item
设  是 A 的素理想， 是 A 的理想。如果  不包含于 \(\dim(A_{p}/\mathfrak{a}A_{p}) = -\infty\)，则有 ；如果  包含于 \(\dim(A_{p}/\mathfrak{a}A_{p}) = \text{codim}(V(p), V(a))\)，则有 。特别地，如果  是 A 的素理想，则有 \(\dim(A_{\mathfrak{p}}) = \operatorname{ht}(\mathfrak{p})\)。
\item
如果  是 A 的理想，则有 ，其中  遍历包含 \_ 的 A 的素理想集合。
\end{enumerate}

断言  是 no. 2 中定义 3  的翻译。断言  由序保持同构这一事实得出：该同构把 A 中满足  的素理想集合  映到局部环  的素理想集合（II, § 2, no. 5, 命题 11）。设  是 A 的一个理想；  的不可约闭子集是集合 ，其中  是包含  的 A 的素理想。因此断言  由 no. 2 的注 1 得出。

\(a\)\(a\)\(b\)\(\mathfrak{q} \mapsto \mathfrak{q}(A_{\mathfrak{p}} / aA_{\mathfrak{p}})\)\(\mathfrak{q}\)\(\mathfrak{a} \subset \mathfrak{q} \subset \mathfrak{p}\)\(A_{\mathfrak{p}} / aA_{\mathfrak{p}}\)\(\mathfrak{a}\)\(\mathfrak{a}\)推论。--- 设 \(\mathfrak{p}\) 是 A 的一个素理想，S 是 A 的一个不交于 \(\mathfrak{p}\) 的乘法子集。于是 \(\mathfrak{a}\)\(c\)

\(^{-1}\).

这由命题 7 a) 以及 II, § 2, no. 5 的命题 11 得出。

命题 8. --- 设 A 是一个环。

\begin{enumerate}
\def\labelenumi{\alph{enumi})}
\item
我们有  ，其中 \(\dim(A) = \sup_{\mathfrak{m}} \dim(A_{\mathfrak{m}}) = \sup_{\mathfrak{m}} \operatorname{ht}(\mathfrak{m})\)\(m\)\(A\) 遍历 A 的极大（分别为素）理想集合。
\item
设  是 A 的一个异于 A 的理想，并设  是 A 的一个包含于  中的理想。则有  。特别地，对 A 的每个异于 A 的理想  ，都有不等式 。
\end{enumerate}

第一条断言由 no. 2 的注记 2 及命题 7 b) 得出。第二条断言由 no. 2 的命题 3 a) 及关系  、  得出。

定义 7. --- 如果拓扑空间 Spec(A) 是串联的（no. 2, 定义 4），则称环 A 是串联环。

这意味着，对于 A 的每一对素理想 \((\mathfrak{p}, \mathfrak{q})\)，只要 \(\mathfrak{q} \subset \mathfrak{p}\)，所有以 \(\mathfrak{p}\) 和 \(\mathfrak{q}\) 为端点的 A 的素理想饱和链，都具有长度 \(\operatorname{codim}(\mathrm{V}(\mathfrak{p}), \mathrm{V}(\mathfrak{q})) = \dim(\mathrm{A}_{\mathfrak{p}} / \mathfrak{q}\mathrm{A}_{\mathfrak{p}})\)。

注。--- 2) 串联环的每个商环都是串联的。对于环 A 为串联环，当且仅当对于 A 的每个素理想  ，环 \(\mathfrak{p}\) 是串联的。

\begin{enumerate}
\def\labelenumi{\arabic{enumi})}
\setcounter{enumi}{2}
\item
  根据 no. 2 的命题 7 b) 和命题 4，环 A 是串联的，当且仅当对于 A 的每个素理想三元组，只要满足相应包含关系，就有 \(\dim(A_{p}/qA_{p}) + \dim(A_{q}/rA_{q}) = \dim(A_{p}/rA_{p})\)。如果 A 是整环且有限维，则 A 是串联的，当且仅当有 \(ht(q) + \dim(A_{p}/qA_{p}) = ht(p)\)，对于 A 的每一对素理想，满足 \((p, q)\)\(\mathfrak{q} \subset \mathfrak{p}\)。事实上，此时拓扑空间 Spec(A) 是不可约且有限维的，应用 no. 2 的命题 4 的推论即可。
\item
  设 A 是一个有限维环，且其所有素理想极大链具有相同长度。则 A 是串联的；对于 A 的每个素理想 \_，有 ；对于满足 \_ 的每一对素理想，有  （no. 2, prop. 5）。
\item
我们将在 § 2, no. 4 中看到，每个域上的有限型代数都是串联环。存在并非串联的局部诺特环（p. 83, 练习 16）。
\end{enumerate}

